\IntroductionGenerale


L’introduction générale doit capter l’intérêt du lecteur et poser les bases du mémoire. Elle doit être claire, précise et progressive.

\textbf{Structure recommandée :}
\begin{itemize}
	\item Contexte et problématique
	\item Présentation du domaine d’étude et de son importance.
	\item Identification du problème auquel votre mémoire répond.
\end{itemize}

\textbf{Objectifs du mémoire :}

\begin{itemize}
	\item Explication des objectifs principaux et secondaires.
	\item Justification de l’intérêt du sujet.
\end{itemize}

\textbf{Méthodologie employée}

\begin{itemize}
	\item Présentation des approches utilisées pour traiter le problème.
	\item Outils et technologies employés.
\end{itemize}

\textbf{Organisation du mémoire}

\begin{itemize}
	\item Brève description du contenu des chapitres.
	\item Indication du cheminement logique du document.
\end{itemize}

